% ==============================================================================
% The Grand Daily Tournament — 2026-09-01 Match (Day 2)
% Locked template: EB Garamond + Garamond-Math + STKaiti (v4.1)
% Week 1 (Day 2) · Variable-Upper Integral Limits & Harmonic Oscillator
% ==============================================================================
\documentclass[a4paper,11pt]{article}

\usepackage{fontspec}
\usepackage{amsmath}
\usepackage{unicode-math}
\defaultfontfeatures{Ligatures=TeX}
\setmainfont{EBGaramond}[
  Path           = {c:/texlive/2024/texmf-dist/fonts/opentype/public/ebgaramond/},
  Extension      = .otf,
  UprightFont    = *-Regular,
  ItalicFont     = *-Italic,
  BoldFont       = *-SemiBold,
  BoldItalicFont = *-SemiBoldItalic,
  Ligatures      = TeX,
  Numbers        = OldStyle
]
\setsansfont{LibertinusSans}[
  Path        = {c:/texlive/2024/texmf-dist/fonts/opentype/public/libertinus-fonts/},
  Extension   = .otf,
  UprightFont = *-Regular,
  ItalicFont  = *-Italic,
  BoldFont    = *-Bold,
  Scale       = MatchLowercase
]
\setmathfont{Garamond-Math.otf}[
  Path  = {c:/texlive/2024/texmf-dist/fonts/opentype/public/garamond-math/},
  Scale = MatchLowercase
]

\usepackage{xeCJK}
\xeCJKsetup{
  CJKmath      = false,
  PunctStyle   = kaiming,
  CheckSingle  = true,
  AutoFakeBold = true,
  AutoFakeSlant= false
}
\setCJKmainfont{STKaiti}[
  Path        = {C:/Windows/Fonts/},
  UprightFont = STKAITI.TTF,
  BoldFont    = STKAITI.TTF,
  ItalicFont  = STKAITI.TTF,
  Script      = Default
]
\setCJKsansfont{Microsoft YaHei}
\setCJKmonofont{FangSong}

\usepackage[margin=1.8cm,headheight=14pt,headsep=12pt,footskip=20pt]{geometry}
\usepackage{booktabs,enumitem,xcolor,graphicx,tabularx}

\usepackage[breakable,skins,hooks]{tcolorbox}
\usepackage{tikz}
\usetikzlibrary{calc,arrows.meta}
\usepackage{fancyhdr,lastpage}
\usepackage{microtype}

\definecolor{navy}{HTML}{0F172A}
\definecolor{slate}{HTML}{1E293B}
\definecolor{royal}{HTML}{1E40AF}
\definecolor{mist}{HTML}{64748B}
\definecolor{border}{HTML}{E2E8F0}
\definecolor{paper}{HTML}{F8FAFC}
\definecolor{forest}{HTML}{065F46}
\definecolor{amber}{HTML}{D97706}
\definecolor{wine}{HTML}{991B1B}
\definecolor{ice}{HTML}{EFF6FF}

\linespread{1.12}
\setlength{\parindent}{0pt}
\setlength{\parskip}{4pt plus 1pt minus 1pt}
\setlength{\abovedisplayskip}{6pt plus 2pt minus 2pt}
\setlength{\belowdisplayskip}{6pt plus 2pt minus 2pt}

\pagestyle{fancy}\fancyhf{}
\fancyhead[L]{\textcolor{mist}{\footnotesize\sffamily 2027 Theoretical Physics · 604 Calculus \& 811 Quantum Mechanics}}
\fancyhead[R]{\textcolor{slate}{\footnotesize\sffamily\bfseries Daily Tournament · 2026-09-01}}
\fancyfoot[L]{\textcolor{mist}{\footnotesize\itshape ``Calculate until the very last step.''}}
\fancyfoot[R]{\textcolor{mist}{\footnotesize\sffamily Page \thepage\ of \pageref{LastPage}}}
\renewcommand{\headrulewidth}{0.4pt}
\renewcommand{\headrule}{\color{border}\hrule width\headwidth height\headrulewidth \vskip-\headrulewidth}

\newtcolorbox{briefing}[1]{%
  enhanced, blanker, breakable,
  left=4mm, right=3mm, top=3mm, bottom=3mm,
  borderline west={2.5pt}{0pt}{royal},
  colback=paper,
  fonttitle=\sffamily\bfseries\color{slate},
  title={#1},
  attach title to upper={\par\vspace{2mm}}
}

\newtcolorbox{problem}[2]{%
  enhanced, breakable,
  colback=white, colframe=border, boxrule=0.6pt, arc=2mm,
  left=4mm, right=4mm, top=3mm, bottom=3mm,
  fonttitle=\sffamily\bfseries\color{slate},
  title={\raisebox{0pt}{\color{royal}\rule{3pt}{10pt}}\hspace{4pt}#1\hfill{\normalfont\footnotesize\color{mist}\itshape #2}},
  colbacktitle=paper, coltitle=slate,
  titlerule=0pt, toptitle=2mm, bottomtitle=2mm
}

\newtcolorbox{solution}[1]{%
  enhanced, breakable,
  colback=white, colframe=border, boxrule=0.6pt, arc=2mm,
  left=4mm, right=4mm, top=3mm, bottom=3mm,
  fonttitle=\sffamily\bfseries\color{forest},
  title={\raisebox{0pt}{\color{forest}\rule{3pt}{10pt}}\hspace{4pt}#1},
  colbacktitle=paper, coltitle=forest,
  titlerule=0pt, toptitle=2mm, bottomtitle=2mm
}

\newtcolorbox{debrief}[1]{%
  enhanced, breakable,
  colback=white, colframe=border, boxrule=0.6pt, arc=2mm,
  left=3.5mm, right=3.5mm, top=2.5mm, bottom=2.5mm,
  fonttitle=\sffamily\bfseries\color{slate},
  title={\raisebox{0pt}{\color{amber}\rule{3pt}{10pt}}\hspace{4pt}#1},
  colbacktitle=paper, coltitle=slate,
  titlerule=0pt, toptitle=2mm, bottomtitle=2mm
}

\newtcolorbox{notebox}{%
  enhanced, blanker,
  left=3mm, right=2mm, top=1.5mm, bottom=1.5mm,
  borderline west={2pt}{0pt}{royal!60},
  colback=ice, fontupper=\small
}

\newtcolorbox{warnbox}{%
  enhanced, blanker,
  left=3mm, right=2mm, top=1.5mm, bottom=1.5mm,
  borderline west={2pt}{0pt}{wine},
  colback=wine!4, fontupper=\small
}

\newcommand{\checkbox}{\ensuremath{\mdlgwhtsquare}}
\newcommand{\alerttag}[1]{\textcolor{wine}{\sffamily\bfseries [#1]}}

\begin{document}

% ==============================================================================
% PAGE 1: COVER
% ==============================================================================
\begin{titlepage}
\thispagestyle{empty}
\begin{tikzpicture}[remember picture,overlay]
  \fill[paper] (current page.south west) rectangle (current page.north east);

  \fill[navy] (current page.north west) rectangle ($(current page.north east)+(0,-3.4cm)$);
  \fill[royal] ($(current page.north west)+(0,-3.4cm)$) rectangle ($(current page.north east)+(0,-3.55cm)$);
  \fill[amber] ($(current page.north west)+(0,-3.55cm)$) rectangle ($(current page.north east)+(0,-3.62cm)$);

  \node[anchor=west, text=white] at ($(current page.north west)+(2.2cm,-1.5cm)$) {
    {\fontsize{12}{14}\selectfont\sffamily\bfseries THE GRAND DAILY TOURNAMENT \textperiodcentered\ 2027}
  };
  \node[anchor=west, text=white!50] at ($(current page.north west)+(2.2cm,-2.3cm)$) {
    {\fontsize{8.5}{11}\selectfont\sffamily UCAS \textperiodcentered\ HIAIS \textperiodcentered\ THEORETICAL PHYSICS ADVANCED TRAINING DOSSIER}
  };

  \begin{scope}[shift={($(current page.center)+(0,1.2cm)$)}]
    \draw[line width=0.5pt, border!80, dashed] (0,0) circle (3.3cm);
    \draw[line width=0.7pt, border] (0,0) circle (2.5cm);
    \draw[line width=0.4pt, royal!30] (0,0) circle (1.6cm);
    \draw[line width=1.2pt, slate!60] (-3.8,-2.2) -- (3.8,2.2);
    \draw[line width=1.2pt, slate!60] (-3.8,2.2) -- (3.8,-2.2);
    \draw[line width=0.7pt, royal!65, -{Stealth[length=2.5mm]}] (-4.2,0) -- (4.2,0)
      node[right, font=\footnotesize\sffamily\color{mist}] {$\hat{x},\, \hat{a}$};
    \draw[line width=0.7pt, royal!65, -{Stealth[length=2.5mm]}] (0,-3.6) -- (0,3.6)
      node[above, font=\footnotesize\sffamily\color{mist}] {$\hat{p},\, \hat{a}^\dagger$};
    \draw[line width=1.3pt, royal, smooth, samples=120, domain=-3.14:3.14]
      plot ({\x}, {1.3*exp(-\x*\x/2.0)});
    \draw[line width=1.1pt, amber!85!black, smooth, samples=120, domain=-3.14:3.14]
      plot ({\x}, {1.3*1.414*\x*exp(-\x*\x/2.0)});
    \fill[navy] (0,0) circle (3.5pt);
    \fill[amber] (0,0) circle (2pt);
    \fill[royal] (1.5, 0.43) circle (2.2pt);
    \fill[wine] (-1.5, -0.43) circle (2.2pt);
    \fill[forest] (0, 1.3) circle (2.2pt);
  \end{scope}

  \node[anchor=center] at ($(current page.center)+(0,-3.4cm)$) {
    \begin{minipage}{0.82\textwidth}
      \centering
      {\fontsize{32}{38}\selectfont\bfseries\color{navy} Daily Tournament}\\[5pt]
      {\fontsize{12}{15}\selectfont\sffamily\bfseries\color{royal}%
        GRAND ARENA OF THEORETICAL PHYSICS}\\[12pt]
      {\color{navy}\rule{0.5\textwidth}{1.2pt}}\\[10pt]
      {\fontsize{10.5}{14}\selectfont\color{slate}\itshape
        ``Today's contestants are ready: 604 Variable-Upper Integral $\times$ 811 Harmonic Oscillator''\\[3pt]
        {\small\color{mist}\upshape Closed-book \textperiodcentered\
        Pen-and-paper derivation \textperiodcentered\ No shortcuts}}
    \end{minipage}
  };

  \node[anchor=south] at ($(current page.south)+(0,1.6cm)$) {
    \begin{minipage}{0.84\textwidth}
      \begin{tcolorbox}[
        enhanced, colback=white, colframe=border, boxrule=0.6pt, arc=2mm,
        left=4mm, right=4mm, top=3.5mm, bottom=3.5mm
      ]
        \renewcommand{\arraystretch}{1.3}
        \sffamily\small
        \begin{tabularx}{\linewidth}{@{}p{2.8cm}X@{}}
          \textbf{Date} & 2026-09-01 \quad\textbar\quad Week 1 (Day 2)\\
          \textbf{Modules} & \textbf{604}: Variable-upper limits \& Taylor \quad\textbar\quad \textbf{811}: Harmonic oscillator wave mechanics\\
          \textbf{Error Protocol} & 5-Factor: Knowledge\,\textbar\, Modelling\,\textbar\, Logic\,\textbar\, Calculation\,\textbar\, Time\\
          \textbf{Target} & Standard 5.5--6.5\,h \quad\textbar\quad Minimum Defence 3.5\,h\\
        \end{tabularx}
      \end{tcolorbox}
    \end{minipage}
  };
\end{tikzpicture}
\end{titlepage}

% ==============================================================================
% PAGE 2: §1 TACTICAL BRIEFING & WARM-UP RETRIEVAL
% ==============================================================================
\clearpage

\begin{center}
  {\LARGE\bfseries\color{navy} Theoretical Physics Training \textperiodcentered\ Daily Tournament}\\[4pt]
  {\small\sffamily\color{mist} 2026-09-01 \quad\textbar\quad Week 1 (Day 2) \quad\textbar\quad Phase: Foundation diagnostic}\\[-3pt]
  {\color{navy}\rule{\textwidth}{0.8pt}}\\[-11pt]
  {\color{border}\rule{\textwidth}{0.3pt}}
\end{center}

\vspace{-4pt}

\begin{briefing}{\S 1\quad Tactical Briefing \& Warm-Up Retrieval}
\begin{itemize}[leftmargin=1.5em, itemsep=2pt, topsep=2pt]
  \item \textbf{604 今日靶向}：变上限积分导数与泰勒展开综合应用。熟练运用 Leibniz 变限求导定理及被积函数高阶展开定阶，攻克 M1 积分分子与等价替换分母之极限。
  \item \textbf{811 今日靶向}：一维简谐振子降算符微分表象 $\hat{a}\psi_0=0$ 解一阶微分方程求高斯基态波函数；利用阶梯代数求解偶极跃迁矩阵元 $\langle m|\hat{x}|n\rangle$ 与测不准关系 $\Delta x \Delta p = \hbar/2$。
  \item \textbf{Supporting}：英语精读 2015 Text 1，掌握长难句三刀切脱水与同义改写；政治学习马原辩证法对立统一规律。
\end{itemize}

\begin{notebox}
\textbf{Momentum Defence Line:} Standard push 5.5--6.5\,h \quad\textbar\quad Minimum defence 3.5\,h (must secure: 811 retrieval + M1 + English reading).
\end{notebox}

\vspace{1mm}
\textbf{Warm-Up Retrieval Checklist} (5--10\,min): without notes, write on scratch paper:\\
\checkbox\ $\frac{d}{dx}\int_0^{g(x)} f(t)\,dt = f(g(x))g'(x)$ \quad
\checkbox\ Annihilation condition: $\hat{a}|0\rangle = 0 \implies \hat{a}\psi_0(x)=0$ \quad
\checkbox\ Ladder action: $\hat{a}|n\rangle=\sqrt{n}|n-1\rangle$, $\hat{a}^\dagger|n\rangle=\sqrt{n+1}|n+1\rangle$
\end{briefing}

\vspace{6pt}

\begin{tcolorbox}[
  enhanced, colback=white, colframe=border, boxrule=0.6pt, arc=2mm,
  left=4mm, right=4mm, top=3mm, bottom=3mm,
  fonttitle=\sffamily\bfseries\color{slate},
  title={\raisebox{0pt}{\color{royal}\rule{3pt}{10pt}}\hspace{4pt}Warm-Up Retrieval Workspace \textbar\ 白纸破冰草稿区},
  colbacktitle=paper, coltitle=slate,
  titlerule=0pt, toptitle=2mm, bottomtitle=2mm
]
\textcolor{mist}{\footnotesize\itshape 5--10 分钟闭卷盲写：在下方空白区域手推上述 3 个核心公式/边界条件，快速越过启动能垒。}
\vspace{9.2cm}
\end{tcolorbox}

% ==============================================================================
% PAGE 3: §2 CORE PROBLEM SET
% ==============================================================================
\clearpage

{\Large\sffamily\bfseries\color{navy} \S 2\quad Core Problem Set}\hfill{\small\sffamily\color{mist}\itshape Complete independently under timed conditions}

\vspace{8pt}

\begin{problem}{Problem M1 \textbar\ 604 Calculus: variable-upper limit}{Suggested time: 25\,min}
求极限
\[
  \lim_{x\to 0}\frac{\displaystyle\int_0^x (t-\sin t)\,dt}{x^2 (e^{x^2}-1)},
\]
要求写出被积函数的泰勒展开定阶依据及分母等价无穷小替换步骤，求出最终精确常数数值。
\end{problem}

\vspace{6pt}

\begin{problem}{Problem Q1 \textbar\ 811 Quantum Mechanics: harmonic oscillator wave mechanics}{Suggested time: 35\,min}
质量为 $m$ 的粒子在势场 $V(x)=\frac{1}{2}m\omega^2 x^2$ 中运动。定义无量纲降算符 $\hat{a} = \sqrt{\frac{m\omega}{2\hbar}}\hat{x} + \frac{i}{\sqrt{2m\hbar\omega}}\hat{p}$。
\begin{enumerate}[leftmargin=1.6em,itemsep=2pt,topsep=2pt]
  \item 由基态条件 $\hat{a}\psi_0(x) = 0$，写出坐标表象下的一阶微分方程，并求出归一化的基态波函数 $\psi_0(x)$（已知高斯积分 $\int_{-\infty}^{+\infty} e^{-\alpha x^2}dx = \sqrt{\pi/\alpha}$）。
  \item 利用坐标算符 $\hat{x} = \sqrt{\frac{\hbar}{2m\omega}}(\hat{a}+\hat{a}^\dagger)$，计算偶极矩阵元 $\langle m|\hat{x}|n\rangle$，并指出跃迁选择定则。
  \item 计算基态下的位置方差 $(\Delta x)^2$ 与动量方差 $(\Delta p)^2$，并验证不确定性原理 $\Delta x \Delta p = \frac{\hbar}{2}$。
\end{enumerate}
\end{problem}

\vspace{6pt}

\begin{problem}{Extension \textbar\ Conceptual Analysis}{Optional \textperiodcentered\ Suggested time: 10\,min}
M1：为何可将被积函数展开后直接逐项积分？依据何种定理？Q1：从宇称对称性说明为何 $\langle n|\hat{x}|n\rangle \equiv 0$。主问题独立作答合计不超过 80 分钟。
\end{problem}

% ==============================================================================
% PAGE 4: §3 MODEL SOLUTIONS & COMMENTARY
% ==============================================================================
\clearpage

{\Large\sffamily\bfseries\color{navy} \S 3\quad Model Solutions \& Commentary}

\vspace{1pt}
{\small\color{mist}\itshape\hspace{1pt} Complete your independent attempt on scratch paper before consulting this page.}

\vspace{8pt}

\begin{solution}{M1 \textperiodcentered\ Model Solution \& Commentary}
\textbf{Strategy:}\enspace 被积函数包含相消项 $t - \sin t$，先对其做 Taylor 展开定阶；分母是乘积形式，可用等价无穷小快速化简。

\vspace{1mm}
\textbf{Derivation:}\enspace
当 $t \to 0$ 时，$\sin t = t - \dfrac{t^3}{6} + o(t^3) \implies t - \sin t = \dfrac{t^3}{6} + o(t^3)$。\\
分子定限积分：$\displaystyle\int_0^x (t - \sin t)\,dt = \int_0^x \left(\frac{t^3}{6} + o(t^3)\right)dt = \frac{x^4}{24} + o(x^4)$。\\
分母乘积项等价代换：$x^2 (e^{x^2} - 1) \sim x^2 \cdot x^2 = x^4$ ($x \to 0$)。故极限为：
\[
  \lim_{x\to 0}\frac{\int_0^x (t-\sin t)\,dt}{x^2 (e^{x^2}-1)} = \lim_{x\to 0}\frac{\frac{x^4}{24} + o(x^4)}{x^4} = \frac{1}{24}.
\]

\begin{warnbox}
\alerttag{Common Error}\enspace 被积函数展开到 $t^3$ 阶，定限积分后直接升至 $x^4$ 阶，恰与分母 $x^4$ 同阶相消。切忌在被积函数内部盲目只保留到一次项。
\end{warnbox}
\end{solution}

\vspace{4pt}

\begin{solution}{Q1 \textperiodcentered\ Model Solution \& Commentary}
\textbf{Physical Picture:}\enspace 降算符作用在基态上为零矢；坐标算符仅含一次阶梯算符，跃迁仅在相邻能级发生（$\Delta n = \pm 1$）。

\vspace{1mm}
\textbf{Derivation:}\enspace
(1) 动量算符 $\hat{p} = -i\hbar \dfrac{d}{dx}$ 代入 $\hat{a}\psi_0 = 0$：
$\left(\sqrt{\frac{m\omega}{2\hbar}}x + \frac{\hbar}{\sqrt{2m\hbar\omega}}\frac{d}{dx}\right)\psi_0 = 0 \implies \dfrac{d\psi_0}{dx} = -\dfrac{m\omega}{\hbar}x \psi_0$。\\
分离变量积分得 $\psi_0(x) = A e^{-\frac{m\omega}{2\hbar}x^2}$。由归一化 $\int_{-\infty}^{+\infty} |\psi_0|^2 dx = A^2 \sqrt{\frac{\pi\hbar}{m\omega}} = 1 \implies A = \left(\frac{m\omega}{\pi\hbar}\right)^{1/4}$：
\[
  \psi_0(x) = \left(\frac{m\omega}{\pi\hbar}\right)^{1/4} \exp\left(-\frac{m\omega}{2\hbar}x^2\right).
\]
(2) 利用 $\hat{x} = \sqrt{\frac{\hbar}{2m\omega}}(\hat{a}+\hat{a}^\dagger)$：
$\langle m|\hat{x}|n\rangle = \sqrt{\frac{\hbar}{2m\omega}}\left(\sqrt{n}\delta_{m, n-1} + \sqrt{n+1}\delta_{m, n+1}\right)$。\\
跃迁选择定则为：**$\Delta n = m - n = \pm 1$**（仅当 $m = n \pm 1$ 时矩阵元非零）。\\
(3) 基态 $|0\rangle$ 下：$\langle 0|\hat{x}|0\rangle = 0, \langle 0|\hat{p}|0\rangle = 0$。方差为 $(\Delta x)^2 = \frac{\hbar}{2m\omega}, (\Delta p)^2 = \frac{m\hbar\omega}{2}$。\\
不确定度乘积为：$\Delta x \Delta p = \sqrt{\frac{\hbar}{2m\omega}} \cdot \sqrt{\frac{m\hbar\omega}{2}} = \frac{\hbar}{2}$，精确达到 Heisenberg 理论下限！

\begin{notebox}
\alerttag{Checkpoint}\enspace 归一化系数为 $(m\omega/\pi\hbar)^{1/4}$（四分之一次幂）。选择定则矩阵元仅两项非零，基态波包为最小测不准高斯波包。
\end{notebox}
\end{solution}

% ==============================================================================
% PAGE 5: §4 DAILY DEBRIEF & REFLECTION
% ==============================================================================
\clearpage

{\Large\sffamily\bfseries\color{navy} \S 4\quad Daily Debrief \& Reflection}

\vspace{6pt}

\begin{debrief}{Post-Match Review (2026-09-01 Match Official Debrief)}
\renewcommand{\arraystretch}{1.25}
\sffamily\small
\begin{tabularx}{\linewidth}{@{}p{2.4cm}X@{}}
  \textbf{604 Calculus} & Time spent: \underline{\enspace\textbf{18}\enspace} min \quad\textbar\quad M1 status: \textbf{[PASS]~Independent} \quad\textbar\quad Score: \textbf{50 / 50}\\
  & Error type: \textbf{[OK]~Zero Error} (Knowledge \textperiodcentered\ Logic \textperiodcentered\ Calculation All Flawless)\\
  \midrule
  \textbf{811 QM} & Time spent: \underline{\enspace\textbf{32}\enspace} min \quad\textbar\quad Q1 status: \textbf{[PASS]~Independent} \quad\textbar\quad Score: \textbf{50 / 50}\\
  & Error type: \textbf{[OK]~Zero Error} (Generalized to arbitrary state $|n\rangle$: $\Delta x\Delta p \ge \hbar/2$)\\
  \midrule
  \textbf{Planning} & \textbf{Spaced retrieval in 3 days:} 9月4日晨间盲写谐振子任意态方差与变限积分定阶\\
  & \textbf{Today's assessment:} \textbf{[PASS]~Peak Performance (100/100 Full Marks)} \quad\textbar\quad Total: \textbf{55\,min}\\
\end{tabularx}
\end{debrief}

\vspace{6pt}

\begin{tcolorbox}[
  enhanced, colback=white, colframe=border, boxrule=0.6pt, arc=2mm,
  left=4mm, right=4mm, top=3.5mm, bottom=3.5mm,
  fonttitle=\sffamily\bfseries\color{slate},
  title={\raisebox{0pt}{\color{forest}\rule{3pt}{10pt}}\hspace{4pt}Key Derivation Insights \& Traps \textbar\ 突破口与避坑沉淀},
  colbacktitle=paper, coltitle=slate,
  titlerule=0pt, toptitle=2mm, bottomtitle=2mm
]
\sffamily\small
\textbf{Core Breakthrough (今日核心突破口与关键一步)}:\\[3pt]
{\color{slate}
1. \textbf{M1 妙解}：先求原函数 $\int_0^x (t-\sin t)dt = \frac{1}{2}x^2+\cos x-1$，再对 $\cos x$ 展开四阶相消，分母代换一步直达 $\frac{1}{24}$！\\
2. \textbf{Q1 泛化}：由降算符微分表象 $\hat{a}\psi_0=0$ 求得基态波函数 $(m\omega/\pi\hbar)^{1/4}e^{-\frac{m\omega}{2\hbar}x^2}$，并将方差成功推广至任意第 $n$ 激发态 $(\Delta x)_n(\Delta p)_n = \frac{\hbar}{2}(2n+1) \ge \frac{\hbar}{2}$！
}\\[8pt]
\textbf{Fatal Trap \& Common Error (致命陷阱与避坑警示)}:\\[3pt]
{\color{wine}
1. 变上限积分被积函数若展开到 $t^k$ 阶，定限积分后直接升至 $x^{k+1}$ 阶，阶数必须与分母严格对齐。\\
2. 宇称不变性要求电偶极矩阵元 $\langle m|\hat{x}|n\rangle$ 仅在相反宇称态间发生（$\Delta n = \pm 1$），对角元 $\langle n|\hat{x}|n\rangle \equiv 0$。
}\\[8pt]
\textbf{Day +3 Action Item (3 天后间隔重做提取的具体题号与断点)}:\\[3pt]
{\color{royal}
9月4日（周五）晨间 5 分钟白纸盲写：变上限积分定阶与谐振子任意态位置动量方差推导断点。
}
\end{tcolorbox}

\end{document}
